\section{Linha do tempo do documentário}
\label{sec:linha-do-tempo}

Resumo do \citetitle{video:java-story}, na ordem do filme. Os números entre
parênteses são a minutagem do vídeo (\emph{hh:mm:ss} ou \emph{mm:ss}), para
você pular direto para o trecho. Onde o filme é vago ou arredonda uma data,
a nota \enquote{conferido} dá a data das fontes primárias; e cada fase
aponta onde a apostila trata o assunto.

\subsection{O nascimento do Oak (1990--1994)}
\begin{description}
\item[1990--1991] James Gosling e uma equipe pequena da Sun Microsystems
  começam o \emph{Green Project}, para eletrônicos de consumo e TV
  interativa. A linguagem se chama \textbf{Oak} (01:33--02:24).
\item[1992--1994] O projeto perde as concorrências de TV a cabo e passa
  por um período em que ninguém sabe se ele continua (02:32--03:46).
\end{description}
Os objetivos de projeto que o white paper de 1996 dá para a linguagem
estão na seção~\ref{sec:objetivos-1995}.

\subsection{A virada para a web (1994--1995)}
\begin{description}
\item[1994] Com o navegador Mosaic, a equipe vê na web o lugar da
  linguagem e escreve o navegador HotJava para levar interatividade às
  páginas (04:02--04:27).
\item[1995] O Java é lançado com a promessa \termo{Write Once, Run
  Anywhere}{\enquote{Escreva uma vez, rode em qualquer lugar}: o mesmo
  \texttt{.class} roda em qualquer sistema que tenha uma JVM.}. O acordo
  com a Netscape para pôr Java no navegador é o que populariza a linguagem
  (06:34--06:57).
\end{description}
O \enquote{rode em qualquer lugar} é a JVM com o bytecode como contrato:
o capítulo~\ref{cap:lendo-bytecode} mostra esse contrato por dentro, e
a seção~\ref{sec:verificacao} conta por que a JVM verifica o bytecode
antes de executá-lo.

\subsection{Conflitos, padronização e ecossistema (1996--2006)}
\begin{description}
\item[1996--1998] A \enquote{guerra do Java} com a Microsoft, que acaba
  nos tribunais porque a implementação para Windows quebrava a
  compatibilidade (18:20--22:45). \emph{Conferido:} a Sun processou a
  Microsoft em outubro de 1997, e o acordo saiu em janeiro de 2001
  (\citetitle{doc:microsoft-sun-acordo-2001}).
\item[1999] O filme fala da criação do \textbf{Java Community Process}
  (JCP), que formaliza a evolução da linguagem com a indústria
  (24:45--26:20). \emph{Conferido:} o JCP foi criado em 1998
  (\citetitle{doc:jcp}); 1999 é o ano em que a Sun lançou o J2EE, a
  plataforma corporativa, em dezembro (\citetitle{doc:sun-j2ee-1999}).
\item[2000--2006] O Java domina o mercado corporativo com o J2EE. Surgem o
  Apache Tomcat (28:59), o Spring e o Hibernate (33:15--35:10).
  \emph{Conferido:} o Tomcat foi doado à Apache em 1999
  (\citetitle{doc:tomcat-heritage}), o Hibernate é de 2001
  (\citetitle{doc:hibernate-sourceforge-2002}) e o Spring~1.0 é de março
  de 2004 (\citetitle{artigo:spring-1-0-final}).
\item[2006] A Sun anuncia que vai abrir o código do Java: nasce o
  \textbf{OpenJDK} (36:07--38:31). \emph{Conferido:} o anúncio, com o
  HotSpot e o \texttt{javac} sob GPL, é de novembro de 2006
  (\citetitle{artigo:reinhold-one-giant-leap}); um JDK compilável por
  completo saiu sob GPL em maio de 2007
  (\citetitle{artigo:reinhold-liftoff}).
\end{description}
Três decisões de concorrência do capítulo~\ref{cap:por-que} são dessa
fase: as APIs de I/O escalável da JSR~51, no Java~1.4 (2002;
seção~\ref{sec:muitos-requests}), e, no Java~5 (2004), o modelo de
memória da JSR~133 (seção~\ref{sec:jmm-historia}) e o
\texttt{java.util.concurrent} da JSR~166 (seção~\ref{sec:juc}).

\paragraph{Por que o J2EE quase fracassou.}
O filme dá três motivos (31:46--33:16):
\begin{itemize}
\item \textbf{Complexidade.} O modelo de programação misturava Java e XML
  e era descrito como doloroso. Muita gente sentia que as especificações
  serviam mais para impor a tecnologia dos fornecedores do que para
  facilitar aplicações simples (31:48--32:02).
\item \textbf{Herança do CORBA.} Vários arquitetos do J2EE vinham do
  CORBA, uma arquitetura anterior de objetos distribuídos, e desenharam a
  plataforma para especialistas em sistemas distribuídos, não para o
  desenvolvedor comum (32:05--32:44).
\item \textbf{Projetos fracassando.} Tantos projetos corporativos falhavam
  que houve um risco real de o Java no servidor acabar (33:06--33:16).
\end{itemize}
O Spring e o Hibernate, de código aberto e mais pragmáticos, simplificaram
o desenvolvimento e mantiveram o Java relevante nas empresas
(33:15--35:10). A lição vale para a Parte~III: o Quarkus é a mesma
pergunta feita de novo, \enquote{quanto disso precisa acontecer em
runtime?}, respondida com trabalho em build time
(capítulo~\ref{cap:build-time}) e com uma arquitetura em que o framework
fica na borda (seção~\ref{sec:circulos-quarkus}).

\subsection{A era Oracle e a modernização (2010 em diante)}
\begin{description}
\item[2010] A Oracle compra a Sun. A comunidade se preocupa, mas a gestão
  traz foco em estabilidade e evolução do ecossistema (42:28--46:52).
  \emph{Conferido:} a compra foi concluída em janeiro de 2010
  (\citetitle{doc:oracle-sun-2010}).
\item[2014] O \textbf{Java~8}, com lambdas, é tratado como um
  renascimento que muda a imagem da linguagem (48:07--49:31). O
  \texttt{invokedynamic} que os lambdas usam está na
  seção~\ref{sec:invokedynamic}.
\item[2017 em diante] Lançamentos a cada seis meses e inovação em
  desempenho e concorrência, como o \textbf{Project Loom}
  (52:23--01:00:00), mantendo o Java relevante para nuvem e IA
  (01:08:44--01:11:29). \emph{Conferido:} Mark Reinhold propôs o ciclo de
  seis meses em setembro de 2017 (\citetitle{artigo:forward-faster}); no
  modelo de seis meses, o lançamento de março de 2018 é o JDK~10
  (\citetitle{doc:jep322}).
\end{description}

\paragraph{Project Loom e virtual threads.}
O filme dedica um trecho ao Loom (00:59:52--01:03:30):
\begin{itemize}
\item \textbf{O runtime troca, a API fica.} O Java mudou por dentro como o
  runtime gerencia threads. Você continua usando \texttt{Thread} e
  \texttt{ExecutorService}; quem decide se a tarefa roda numa
  \emph{platform thread} (do sistema operacional) ou numa \emph{virtual
  thread} (gerenciada pela JVM) é quem cria a thread (01:02:16--01:02:40).
\item \textbf{Fim da obrigação do reativo.} Antes, para escalar, era
  preciso programação reativa, descrita no filme como barroca e difícil
  de depurar. Com virtual threads, o código volta a ser imperativo e
  sequencial e continua escalando (01:02:43--01:03:00).
\item \textbf{Programas produtivos.} Josh Long resume: o Java~8 deixou o
  \emph{programador} mais produtivo; o Loom deixa o \emph{programa} mais
  produtivo e escalável (01:03:10--01:03:30).
\end{itemize}
\emph{Conferido:} as virtual threads entraram como preview no JDK~19
(2022) e ficaram finais no JDK~21, em setembro de 2023
(\citetitle{doc:jep444}). O bytecode e as regras de uso estão na
seção~\ref{sec:virtual-threads}; a escolha entre virtual threads e reativo
no Quarkus, na seção~\ref{sec:vt-vs-reativo}; e o caminho das green
threads às virtual threads, na seção~\ref{sec:threads-ate-vt}.
